\section{Kualitas Kebijakan}

Hipotesis H1 menyatakan CoDA-EDA menghasilkan kualitas kebijakan yang lebih tinggi dibanding keluarga \textit{compact} univariat, dengan keunggulan yang menguat seiring dimensi ruang solusi, serta setara iBOA/BOA pada anggaran aturan memadai.
Kualitas diukur melalui \textit{F-score}, cakupan (\textit{coverage}), dan kompleksitas struktural (WSC), dengan uji Friedman diikuti post-hoc Wilcoxon berkoreksi Holm-Bonferroni atas ketujuh kondisi algoritma.
Gambaran menyeluruh rerata \textit{F-score} pada kelima dataset disajikan pada Tabel \ref{tab:rerata-fscore}.

\begin{longtable}{|l|p{2cm}|p{2cm}|p{2cm}|p{2cm}|p{2cm}|}
\caption{Rerata \textit{F-score} Ketujuh Algoritma pada Lima Dataset ABAC Lab}
\label{tab:rerata-fscore} \\
\hline
\textbf{Algoritma} & \textbf{workforce (436)} & \textbf{e-document (1.907)} & \textbf{healthcare (46)} & \textbf{project-mgmt (51)} & \textbf{university (54)} \\
\hline
\endfirsthead

\multicolumn{6}{c}{{\bfseries \tablename\ \thetable{} -- lanjutan dari halaman sebelumnya}} \\
\hline
\textbf{Algoritma} & \textbf{workforce (436)} & \textbf{e-document (1.907)} & \textbf{healthcare (46)} & \textbf{project-mgmt (51)} & \textbf{university (54)} \\
\hline
\endhead

\hline
\multicolumn{6}{r}{Lanjut ke halaman berikutnya} \\
\endfoot

\hline
\endlastfoot

% ==================== DATA BARIS ====================
CoDA-EDA & 0{,}889 & 0{,}452 & 0{,}779 & 0{,}898 & 0{,}830 \\
\hline
compact-GA & 0{,}900 & 0{,}329 & 0{,}765 & 0{,}883 & 0{,}822 \\
\hline
compact-PSO & 0{,}893 & 0{,}260 & 0{,}801 & 0{,}902 & 0{,}829 \\
\hline
iBOA & 0{,}888 & 0{,}452 & 0{,}802 & 0{,}909 & 0{,}824 \\
\hline
BOA & 0{,}832 & 0{,}633 & 0{,}835 & 0{,}927 & 0{,}831 \\
\hline
ACO & 0{,}870 & 0{,}661 & 0{,}761 & 0{,}886 & 0{,}835 \\
\hline
UBK & 0{,}853 & 0{,}630 & 0{,}782 & 0{,}853 & 0{,}869 \\
\hline

\end{longtable}

Pada dua dataset utama, uji Friedman signifikan (\textit{workforce} $p = 2{,}0 \times 10^{-14}$; \textit{e-document} $p = 1{,}6 \times 10^{-42}$). Di \textit{workforce} ($D = 436$) CoDA-EDA ($0{,}889$) mengungguli secara signifikan BOA ($r = +0{,}56$) dan UBK ($r = +0{,}33$), sementara setara dengan \textit{compact-GA}, \textit{compact-PSO}, iBOA, dan ACO. 
Ia berada di papan atas tanpa dominan. 
Di \textit{e-document} ($D = 1.907$), pola berbalik: CoDA-EDA ($0{,}452$) setara iBOA ($0{,}452$), mengungguli \textit{compact-GA} ($r = +0{,}42$) dan \textit{compact-PSO} ($r = +0{,}60$), namun kalah signifikan terhadap BOA, ACO, dan UBK dengan \textit{effect size} besar ($r = -0{,}87$ hingga $-0{,}91$). 
Inilah defisit dimensi-besar yang menjadi ciri utama H1.

Ketiga dataset kecil berperan sebagai konteks pelengkap dan hasilnya menegaskan prediksi pra-eksperimen.
Karena \textit{permit} uniknya sedikit (14-90) sehingga duplikasi entri tinggi menjenuhkan kualitas, pembeda antar-algoritma tumpul.
CoDA-EDA setara dengan mayoritas \textit{baseline} pada ketiganya.
Ia hanya kalah tipis-signifikan dari UBK pada \textit{university} (-0,039), sekaligus mengungguli UBK pada \textit{project-management} (+0,045).
Karena duplikasi menjenuhkan ruang solusi dan menumpulkan pembeda antar-algoritma, hasil pada dataset kecil tidak dijadikan tumpuan vonis H1.

Perincian \textit{e-document} (Tabel \ref{tab:rincian-h1-edocument}) memperlihatkan sumber mekanistik kekalahan pada dimensi besar.
Kelemahan CoDA-EDA bukan pada presisi melainkan pada cakupan/\textit{recall}.
Dengan batas empat aturan, ACO, UBK, dan BOA menemukan aturan luas yang mencakup 47-62\% permintaan, sedangkan CoDA-EDA dan iBOA menghasilkan aturan berpresisi tinggi namun sempit yang hanya mencakup sekitar 25\%.

\begin{longtable}{|>{\raggedright\arraybackslash}p{2cm}|>{\raggedright\arraybackslash}p{1.3cm}|>{\raggedright\arraybackslash}p{1.3cm}|>{\raggedright\arraybackslash}p{1.3cm}|>{\raggedright\arraybackslash}p{1.4cm}|>{\raggedright\arraybackslash}p{1cm}|>{\raggedright\arraybackslash}p{3cm}|}
\caption{Rincian Kualitas H1 pada e-document (D = 1.907; rerata lintas korelasi; post-hoc vs CoDA-EDA)}
\label{tab:rincian-h1-edocument} \\
\hline
\textbf{Algoritma} & \textbf{F-score} & \textbf{$\Delta$F vs CoDA} & \textbf{Cakup-an} & \textbf{$\Delta$ cakup-an vs CoDA} & \textbf{WSC} & \textbf{Posisi terhadap CoDA-EDA (post-hoc Holm)} \\
\hline
\endfirsthead

\multicolumn{7}{c}{{\bfseries \tablename\ \thetable{} -- lanjutan dari halaman sebelumnya}} \\
\hline
\textbf{Algoritma} & \textbf{F-score} & \textbf{$\Delta$F vs CoDA} & \textbf{Cakup-an} & \textbf{$\Delta$ cakup-an vs CoDA} & \textbf{WSC} & \textbf{Posisi terhadap CoDA-EDA (post-hoc Holm)} \\
\hline
\endhead

\hline
\multicolumn{7}{r}{Lanjut ke halaman berikutnya} \\
\endfoot

\hline
\endlastfoot

% ==================== DATA BARIS ====================
ACO & 0{,}661 & $+$0{,}209 & 0{,}572 & $+$0{,}324 & 16{,}9 & Unggul signifikan ($r = -0{,}91$) \\
\hline
BOA & 0{,}633 & $+$0{,}181 & 0{,}624 & $+$0{,}376 & 11{,}2 & Unggul signifikan ($r = -0{,}87$) \\
\hline
UBK & 0{,}630 & $+$0{,}178 & 0{,}468 & $+$0{,}220 & 6{,}2 & Unggul signifikan ($r = -0{,}91$) \\
\hline
iBOA & 0{,}452 & 0{,}000 & 0{,}246 & $-0{,}002$ & 13{,}8 & Setara (n.s.) \\
\hline
CoDA-EDA & 0{,}452 & --- & 0{,}248 & --- & 11{,}8 & --- (acuan) \\
\hline
compact-GA & 0{,}329 & $-0{,}123$ & 0{,}156 & $-0{,}092$ & 10{,}2 & CoDA-EDA unggul ($r = +0{,}42$) \\
\hline
compact-PSO & 0{,}260 & $-0{,}192$ & 0{,}116 & $-0{,}132$ & 11{,}0 & CoDA-EDA unggul ($r = +0{,}60$) \\
\hline

\end{longtable}

Karena kelemahan bersumber dari cakupan, dua tuas hiperparameter diuji pada \textit{e-document} untuk memeriksa apakah defisit dapat diperbaiki dan diterapkan seragam pada seluruh algoritma.
Tuas bobot \textit{false positive} $w_\text{fp}$ tidak berpengaruh (cakupan CoDA-EDA 0,186 menjadi 0,187), menandakan persoalannya bukan pada pembobotan presisi.
Sebaliknya, tuas anggaran jumlah aturan berdampak besar (Tabel \ref{tab:efek-aturan}).
Dengan menaikkan dari empat ke delapan lalu dua belas aturan melonjakkan \textit{F-score} CoDA-EDA dari 0,362 ke 0,571 dan menyempitkan jaraknya terhadap ACO dari -0,301 menjadi -0,100.

\begin{longtable}{|>{\raggedright\arraybackslash}p{2cm}|>{\raggedright\arraybackslash}p{1.2cm}|>{\raggedright\arraybackslash}p{1.3cm}|>{\raggedright\arraybackslash}p{2cm}|>{\raggedright\arraybackslash}p{1cm}|>{\raggedright\arraybackslash}p{1cm}|>{\raggedright\arraybackslash}p{3.5cm}|}
\caption{Efek Anggaran Jumlah Aturan terhadap Kualitas CoDA-EDA pada \textit{e-document}}
\label{tab:efek-aturan} \\
\hline
\textbf{Jumlah aturan} & \textbf{F-score} & \textbf{Cakup-an} & \textbf{$\Delta$ cakupan vs 4-aturan} & \textbf{WSC} & \textbf{Waktu} & \textbf{Posisi terhadap baseline (post-hoc Holm)} \\
\hline
\endfirsthead

\multicolumn{7}{c}{{\bfseries \tablename\ \thetable{} -- lanjutan dari halaman sebelumnya}} \\
\hline
\textbf{Jumlah aturan} & \textbf{F-score} & \textbf{Cakup-an} & \textbf{$\Delta$ cakupan vs 4-aturan} & \textbf{WSC} & \textbf{Waktu} & \textbf{Posisi terhadap baseline (post-hoc Holm)} \\
\hline
\endhead

\hline
\multicolumn{7}{r}{Lanjut ke halaman berikutnya} \\
\endfoot

\hline
\endlastfoot

% ==================== DATA BARIS ====================
4 (baku Tabel \ref{tab:rincian-h1-edocument}) & 0{,}362 & 0{,}186 & --- & 9{,}6 & 97 s & Kalah dari ACO, UBK, BOA; setara iBOA \\
\hline
8 & 0{,}571 & 0{,}391 & $+$0{,}205 & 18{,}8 & 117 s & Kalah dari ACO, UBK; setara BOA \& iBOA \\
\hline
12 & 0{,}568 & 0{,}428 & $+$0{,}242 & 26{,}8 & 136 s & Kalah dari ACO, UBK; setara BOA \& iBOA \\
\hline

\end{longtable}

Pada dua belas aturan, post-hoc menempatkan CoDA-EDA setara iBOA ($p_{\text{adj}} = 0{,}52$), padahal pada empat aturan BOA mengungguli telak, sehingga CoDA-EDA hanya kalah moderat terhadap ACO dan UBK.
Perbaikan ini memiliki harga pada waktu (97 menjadi 136 detik) dan kompleksitas (WSC 9,6 menjadi 26,8)

Vonis H1 adalah terdukung secara bersyarat. 
Hasil eksperimen tidak mendukung dugaan awal bahwa keunggulan CoDA-EDA meningkat secara konsisten seiring bertambahnya tingkat korelasi, dan juga tidak menunjukkan dominasi CoDA-EDA terhadap seluruh kelas algoritma pembanding. 
Sebaliknya, pola yang paling jelas adalah bahwa manfaat CoDA-EDA terhadap keluarga \textit{compact} univariat bergantung pada dimensi ruang solusi. 
Pada \textit{workforce} ($D=436$), CoDA-EDA memperoleh F-score 0,889, yang secara praktis setara dengan \textit{compact-GA} (0,900) dan \textit{compact-PSO} (0,893). 
Sebaliknya, pada \textit{e-document} ($D=1.907$), CoDA-EDA mempertahankan F-score 0,452, sedangkan \textit{compact-GA} dan \textit{compact-PSO} masing-masing turun menjadi 0,329 dan 0,260. 
Dengan demikian, selisih CoDA-EDA terhadap rerata kedua algoritma \textit{compact} meningkat dari sekitar $-0{,}008$ pada \textit{workforce} menjadi sekitar $+0{,}158$ pada \textit{e-document}.

Analisis menurut tingkat korelasi memperkuat interpretasi tersebut. 
Pada \textit{e-document}, selisih F-score CoDA-EDA terhadap rerata \textit{compact-GA} dan \textit{compact-PSO} tetap berada pada kisaran sekitar $+0{,}15$ hingga $+0{,}17$ pada korelasi rendah, sedang, maupun tinggi. 
Selisih yang relatif stabil tersebut menunjukkan bahwa keunggulan CoDA-EDA pada dataset ini tidak semakin besar hanya karena tingkat korelasi dinaikkan. 
Faktor pembeda yang lebih konsisten adalah besarnya dimensi ruang solusi. 
Pada ruang berdimensi tinggi, pencuplikan berdasarkan probabilitas marginal secara independen memiliki peluang lebih besar menghasilkan kombinasi pasangan atribut-nilai yang kurang sesuai dengan struktur data. 
CoDA-EDA mengurangi keterbatasan tersebut dengan memanfaatkan dependensi orde satu yang dipelajari melalui \textit{mutual information} dan direpresentasikan dalam pohon Chow--Liu sebagai informasi tambahan ketika membangkitkan kandidat aturan.

Temuan ini tidak berarti bahwa CoDA-EDA menjadi algoritma dengan kualitas penambangan statis terbaik pada dimensi tinggi. 
Pada \textit{e-document}, CoDA-EDA tetap setara dengan iBOA, tetapi tertinggal dari BOA, ACO, dan UBK pada konfigurasi empat aturan. 
Perincian hasil menunjukkan bahwa kelemahan tersebut terutama berasal dari cakupan aturan: CoDA-EDA menghasilkan aturan dengan presisi tinggi tetapi cakupan relatif sempit, sedangkan BOA, ACO, dan UBK mampu menemukan aturan yang mencakup bagian log lebih besar. 
Ketika anggaran aturan ditingkatkan menjadi delapan dan dua belas aturan, jarak kualitas terhadap sebagian algoritma tersebut menyempit, tetapi dibayar dengan peningkatan waktu komputasi dan kompleksitas struktural. 
Oleh sebab itu, dukungan H1 dibatasi pada keunggulan pemodelan dependensi terhadap pendekatan \textit{compact} univariat pada ruang solusi berdimensi tinggi, bukan pada dominasi terhadap seluruh metaheuristik.

Pola pada validasi eksternal juga konsisten dengan batas klaim tersebut. 
Dataset \textit{Smart Building IoT} memiliki dimensi efektif sekitar $D=428$, yang sebanding dengan \textit{workforce}, dan tidak menggunakan korelasi hasil injeksi generator penelitian. 
Pada dataset ini CoDA-EDA memperoleh F-score 0,938, setara secara praktis dengan iBOA sebesar 0,943 dan sedikit di bawah \textit{compact-GA} dan \textit{compact-PSO} yang masing-masing memperoleh 0,961 dan 0,964. 
Hasil tersebut tidak menunjukkan keunggulan CoDA-EDA terhadap keluarga \textit{compact} pada dimensi menengah, sehingga konsisten dengan temuan ABAC Lab bahwa manfaat pemodelan dependensi baru terlihat jelas pada kondisi berdimensi sangat tinggi.

Secara keseluruhan, H1 menunjukkan adanya pertukaran yang terkait langsung dengan keputusan rancangan CoDA-EDA. 
Representasi probabilistik yang diperkaya dengan pohon Chow--Liu memberikan informasi struktural yang tidak tersedia pada model \textit{compact} univariat dan memberikan keuntungan nyata ketika ruang solusi menjadi sangat besar. 
Namun, penggunaan model ringkas tanpa populasi eksplisit juga membatasi keragaman eksplorasi dibanding algoritma berpopulasi seperti ACO dan UBK. 
Dengan demikian, posisi CoDA-EDA bukan sebagai penambang kebijakan statis yang selalu menghasilkan F-score tertinggi, melainkan sebagai algoritma \textit{dependency-aware} yang mempertahankan kualitas kompetitif, menunjukkan keunggulan terhadap pendekatan \textit{compact} univariat pada dimensi tinggi, dan sekaligus menyediakan karakteristik efisiensi memori serta pembaruan yang menjadi fokus utama penelitian ini.

Sebagai ilustrasi keluaran konkret proses penambangan, Tabel \ref{tab:contoh-kebijakan} menyajikan cuplikan aturan yang dihasilkan CoDA-EDA pada \textit{workforce} dan \textit{Smart Building IoT} beserta pembandingnya dari ACO.
Semantik pencocokan mengikuti bahasa ABAC Lab yang menyatakan bahwa beberapa nilai pada atribut yang sama bersifat disjungtif dan antar-atribut bersifat konjungtif.

\begin{longtable}{|>{\raggedright\arraybackslash}p{2.5cm}|>{\raggedright\arraybackslash}p{2.5cm}|>{\raggedright\arraybackslash}p{5cm}|>{\raggedright\arraybackslash}p{2cm}|}
\caption{Contoh Kebijakan Hasil Mining (cuplikan aturan terbaca)}
\label{tab:contoh-kebijakan} \\
\hline
\textbf{Sumber / algoritma} & \textbf{Metrik} & \textbf{Contoh aturan (cuplikan)} & \textbf{Cakupan aturan} \\
\hline
\endfirsthead

\multicolumn{4}{c}{{\bfseries \tablename\ \thetable{} -- lanjutan dari halaman sebelumnya}} \\
\hline
\textbf{Sumber / algoritma} & \textbf{Metrik} & \textbf{Contoh aturan (cuplikan)} & \textbf{Cakupan aturan} \\
\hline
\endhead

\hline
\multicolumn{4}{r}{Lanjut ke halaman berikutnya} \\
\endfoot

\hline
\endlastfoot

% ==================== DATA BARIS ====================
workforce / CoDA-EDA & F1 0{,}885; WSC 20; 5 aturan & IF contractStatus=active AND action=view THEN permit & TP=1.291; FP=339 \\
\hline
workforce / CoDA-EDA & (aturan berpresisi) & IF provider in \{externalHelpdeskSupplier, telco\} AND tenant=telco AND associatedContract in \{contract001,003,005\} AND action=modify THEN permit & TP=402; FP=3 \\
\hline
workforce / ACO & F1 0{,}820; WSC 19; 5 aturan & IF tenantType=primary AND action in \{modify,view\} THEN permit & TP=2.879; FP=1.297 (luas, presisi rendah) \\
\hline
Smart Building IoT / CoDA-EDA & F1 0{,}998; WSC 276; 5 aturan & IF type in \{TV, Security Cameras, Smart locks\} AND action in \{arm, control\} THEN permit & TP=3.417; FP=0 \\
\hline

\end{longtable}

Cuplikan ini memperlihatkan sifat keluaran yang berbeda antar-algoritma.
CoDA-EDA cenderung menghasilkan aturan berpresisi tinggi (mis. FP=3 dari 405 pada aturan \textit{modify workforce}, FP=0 pada IoT), sedangkan ACO cenderung menghasilkan aturan lebih luas bercakupan tinggi namun berpresisi lebih rendah (FP=1.297).
Hal tersebut konsisten dengan analisis mekanistik cakupan-presisi pada Tabel \ref{tab:rincian-h1-edocument}.
Perlu dicatat pula bahwa pada \textit{Smart Building IoT} sebagian aturan CoDA-EDA membengkak WSC-nya (total 276) karena mengenumerasi nilai atribut bersifat pengenal (uname/rname) sebagai disjungsi panjang.